\question[2]

a. Was sagt die Fano-Bedinung aus?

b. Unten sind Codewörter für einige Buchstaben gegeben. 
Ist die Fano-Bedingung erfüllt? Finde ein Beispiel für eine Codierung, die nicht eindeutig dekodiert werden kann.

a = '00' ~~ b = '10' ~~ c = '11' ~~ d = '01' ~~ f = '101'
 
\begin{solutionbox}{4cm}
a. Die Fanobedingung besagt, dass kein Codewort der Anfang eines anderen Codewortes sein darf. Wenn die Fanobedingung erfüllt ist, lassen sich Codierungen eindeutig decodieren.
 
b. Das Codewort für b ist der Anfang des Codeworts für f. Die Fano-Bedingung ist nicht erfüllt.
Beispiel für eine nicht eindeutige Codierung: 101101 kann mit bcd oder ff dekodiert werden.
\end{solutionbox}